\chapter{Framework for a Minimal FDTD Solver for Linearized MHD of the Polomac}

\section{Problem formulation}
We want to be able to solve for the independent time dependent properties of the MHD equations in \autoref{subsec: The ideal MHD equations},
which are $\varrho$, $\vec{V}$, $\vec{B}$ and $P$.  


If we have $\vec{B}$ and $\vec{V}$ we will be able to solve for $\vec{E}$ and $\vec{J}$ at any time via \eqref{eq: MHD ohm's law} and \eqref{eq: MHD ampere-maxwell}, respectively.


\subsection{Equations to be solved}
As in \autoref{subsec: theory of alfven waves} we eliminate $\vec{E}$ and $\vec{J}$.
Additionally we take the mass continuity equation \eqref{eq: MHD mass continuity eq} and the closure equation \eqref{eq: MHD closure eq}
to arrive at a system of four first-order hyperbolic partial differential equations
\begin{align*}
    \frac{\partial \varrho}{\partial t} + \nabla \cdot (\varrho \vec{V}) &= 0 \\
    \varrho \frac{d\vec{V}}{dt} &= -\nabla P + \frac{1}{\mu_0}(\nabla \times \vec{B}) \times \vec{B} \\
    \frac{\partial \vec{B}}{\partial t} &= \nabla \times (\vec{V} \times \vec{B}) \\
    \frac{d}{dt}\!\left(P \varrho^{-\gamma}\right) &= 0
\end{align*}
and as an additional constraint Gauss' law for magnetism \eqref{eq: MHD gauss law for magnetism},
which is not evolved but must be preserved
\begin{align*}
    \nabla \cdot \vec{B} &= 0.
\end{align*}



\section{Linearization and discretization}

As in \autoref{subsec: theory of alfven waves}, we linearize around a uniform static equilibrium
\begin{align*}
    \vec{B}(\vec{x}, t) &= \vec{B}_0 + \vec{B}_1(\vec{x}, t), \\
    \vec{V}(\vec{x}, t) &= \vec{V}_1(\vec{x}, t), \\
    \varrho(\vec{x}, t) &= \varrho_0 + \varrho_1(\vec{x}, t), \\
    P(\vec{x}, t)       &= P_0 + P_1(\vec{x}, t),
\end{align*}
where $\vec{B}_0$, $\varrho_0$, $P_0$ are constants, $\vec{V}_0 = 0$, and all perturbations
satisfy $|Q_1/Q_0| \ll 1$. Substituting and retaining only first-order terms yields
the linearized equations
\begin{align}
    \frac{\partial \varrho_1}{\partial t} &= -\varrho_0\, \nabla \cdot \vec{V}_1,
        \label{eq: lin continuity} \\
    \varrho_0\, \frac{\partial \vec{V}_1}{\partial t} &= -\nabla P_1
        + \frac{1}{\mu_0}\bigl(\nabla \times \vec{B}_1\bigr) \times \vec{B}_0,
        \label{eq: lin momentum} \\
    \frac{\partial \vec{B}_1}{\partial t} &= \nabla \times \bigl(\vec{V}_1 \times \vec{B}_0\bigr),
        \label{eq: lin induction} \\
    \frac{\partial P_1}{\partial t} &= -\gamma P_0\, \nabla \cdot \vec{V}_1.
        \label{eq: lin pressure}
\end{align}
Equations \eqref{eq: lin continuity} and \eqref{eq: lin pressure} together imply
$P_1 = c_s^2 \varrho_1$, where $c_s = \sqrt{\gamma P_0 / \varrho_0}$ is the adiabatic
sound speed, so they carry the same information. The system is therefore closed by
\eqref{eq: lin momentum} and \eqref{eq: lin induction}, with $P_1$ supplied
by \eqref{eq: lin pressure}. The constraint $\nabla \cdot \vec{B}_1 = 0$ is
automatically preserved by \eqref{eq: lin induction} and serves as a numerical
consistency check. The unknowns to be time-stepped are $\vec{V}_1$, $\vec{B}_1$,
and $P_1$.


\subsection{FDTD}

\subsection{Types of PDEs}
\label{subsec: types of PDEs}

The linearized system \eqref{eq: lin continuity}--\eqref{eq: lin pressure} consists
of first-order, linear, hyperbolic partial differential equations. This can be seen
by noting that all time derivatives appear linearly on the left-hand side and the
right-hand sides contain only spatial derivatives of the unknowns, with no
dissipative (parabolic) or elliptic terms. Hyperbolic systems propagate information
at finite speeds along characteristics; the relevant speeds here are the Alfv\'en
velocity $v_A$ and the magnetosonic speed $v_f = \sqrt{v_A^2 + c_s^2}$, both
derived in \autoref{sec: relevant parameters}. This structure makes the system
well-suited to explicit time-domain methods such as FDTD.

\subsection{Reduction to 2D in cylindrical coordinates}
\label{subsec: 2D cylindrical reduction}

The Polomac geometry is naturally described in cylindrical coordinates
$(r, \varphi, z)$. We adopt the axisymmetric approximation already used in
\autoref{chap: static equilibrium}, setting $\partial/\partial\varphi = 0$.
All perturbation fields are then functions of $(r, z, t)$ only.

We choose the equilibrium field along the $z$-axis,
$\vec{B}_0 = B_0 \hat{e}_z$, which is the simplest geometry consistent with
the poloidal dominance of the Polomac (see \autoref{subsec: polomac geometry})
and with the Alfv\'en wave analysis in \autoref{subsec: theory of alfven waves}.
In this geometry the axisymmetric ($\partial_\varphi = 0$) cylindrical
differential operators \eqref{eq: cylindrical divergence}, \eqref{eq: cylindrical curl}
reduce to
\begin{align}
    \nabla \cdot \vec{F} &= \frac{1}{r}\frac{\partial(r F_r)}{\partial r}
        + \frac{\partial F_z}{\partial z}, \\
    (\nabla \times \vec{F})_r &= -\frac{\partial F_\varphi}{\partial z}, \qquad
    (\nabla \times \vec{F})_\varphi = \frac{\partial F_r}{\partial z}
        - \frac{\partial F_z}{\partial r}, \qquad
    (\nabla \times \vec{F})_z = \frac{1}{r}\frac{\partial(r F_\varphi)}{\partial r}.
\end{align}

We expand the perturbation fields into components
$\vec{V}_1 = V_r\hat{e}_r + V_\varphi\hat{e}_\varphi + V_z\hat{e}_z$ and
$\vec{B}_1 = B_r^{(1)}\hat{e}_r + B_\varphi^{(1)}\hat{e}_\varphi + B_z^{(1)}\hat{e}_z$.
Inserting $\vec{B}_0 = B_0\hat{e}_z$ into the linearized system
\eqref{eq: lin continuity}--\eqref{eq: lin pressure} and evaluating each component
yields the following scalar equations.

\textbf{Pressure} (from \eqref{eq: lin pressure}):
\begin{align}
    \frac{\partial P_1}{\partial t} &= -\gamma P_0
        \left(\frac{1}{r}\frac{\partial(r V_r)}{\partial r}
        + \frac{\partial V_z}{\partial z}\right).
    \label{eq: 2D pressure}
\end{align}

\textbf{Momentum} (from \eqref{eq: lin momentum}),
using $(\nabla\times\vec{B}_1)\times\vec{B}_0
= B_0\,\partial_z\vec{B}_1^\text{pol} - B_0(\nabla\cdot\vec{B}_1^\text{pol})\hat{e}_z$
evaluated component by component:
\begin{align}
    \varrho_0\frac{\partial V_r}{\partial t} &= -\frac{\partial P_1}{\partial r}
        + \frac{B_0}{\mu_0}\left(\frac{\partial B_r^{(1)}}{\partial z}
        - \frac{\partial B_z^{(1)}}{\partial r}\right)
        \cdot(-1) \cdot(-1),
        \nonumber \\
    &= -\frac{\partial P_1}{\partial r}
        + \frac{B_0}{\mu_0}\frac{\partial B_r^{(1)}}{\partial z}
        - \frac{B_0}{\mu_0}\frac{\partial B_z^{(1)}}{\partial r},
    \label{eq: 2D mom r} \\
    \varrho_0\frac{\partial V_\varphi}{\partial t} &=
        \frac{B_0}{\mu_0}\frac{\partial B_\varphi^{(1)}}{\partial z},
    \label{eq: 2D mom phi} \\
    \varrho_0\frac{\partial V_z}{\partial t} &= -\frac{\partial P_1}{\partial z}.
    \label{eq: 2D mom z}
\end{align}

\textbf{Induction} (from \eqref{eq: lin induction}),
using $\nabla\times(\vec{V}_1\times B_0\hat{e}_z)
= B_0\,\partial_z\vec{V}_1 - B_0(\nabla\cdot\vec{V}_1)\hat{e}_z$:
\begin{align}
    \frac{\partial B_r^{(1)}}{\partial t} &=
        B_0\frac{\partial V_r}{\partial z},
    \label{eq: 2D ind r} \\
    \frac{\partial B_\varphi^{(1)}}{\partial t} &=
        B_0\frac{\partial V_\varphi}{\partial z},
    \label{eq: 2D ind phi} \\
    \frac{\partial B_z^{(1)}}{\partial t} &=
        -B_0\left(\frac{1}{r}\frac{\partial(r V_r)}{\partial r}
        + \frac{\partial V_z}{\partial z}\right).
    \label{eq: 2D ind z}
\end{align}

Note that the $\varphi$-components \eqref{eq: 2D mom phi} and \eqref{eq: 2D ind phi}
decouple from the rest and describe a pure shear Alfv\'en wave propagating along
$\hat{e}_z$, consistent with the dispersion relation \eqref{eq: alfven dispersion}.
The remaining components $(V_r, V_z, P_1, B_r^{(1)}, B_z^{(1)})$ describe the
compressional (magnetosonic) branch.

\subsection{FDTD Discretization}
\label{subsec: FDTD}

\subsubsection{Staggered grid}

We discretize the $(r, z)$ domain on a uniform staggered (Yee-type) grid with
spacings $\Delta r$ and $\Delta z$. Scalar quantities $P_1$, $B_z^{(1)}$, $\varrho_1$
and $V_z$ are located at cell centres $(r_i, z_j)$; radial components $V_r$ and
$B_r^{(1)}$ are located at the $r$-faces $(r_{i+1/2}, z_j)$; azimuthal components
$V_\varphi$ and $B_\varphi^{(1)}$ are co-located with $V_r$ and $B_r^{(1)}$.
Time is staggered by half a step: $\vec{B}_1$ is evaluated at integer times $t^n$
and $\vec{V}_1$, $P_1$ at half-integer times $t^{n+1/2}$. This leapfrog arrangement
ensures that the update of $\vec{B}_1$ from the induction equations uses
$\vec{V}_1$ values that are centred in time, giving second-order accuracy in both
space and time without any implicit solve.

\subsubsection{Courant stability condition}

Because the linearized ideal MHD equations are hyperbolic, explicit time-stepping
is subject to the Courant--Friedrichs--Lewy (CFL) stability condition. The fastest
signal speed in the system is the fast magnetosonic speed
\begin{align}
    v_f = \sqrt{v_A^2 + c_s^2},
    \label{eq: fast speed}
\end{align}
where $v_A$ is the Alfv\'en velocity \eqref{eq: alfven velocity} and
$c_s = \sqrt{\gamma P_0/\varrho_0}$ is the adiabatic sound speed.
In two spatial dimensions the CFL condition reads
\begin{align}
    \Delta t \leq \frac{1}{v_f}
        \left(\frac{1}{\Delta r^2} + \frac{1}{\Delta z^2}\right)^{-1/2}.
    \label{eq: CFL}
\end{align}
Using the Polomac target values from \autoref{chap: relevant parameters}
($v_A \approx 5.5\times10^5\ \mathrm{m\,s^{-1}}$,
$T_i \approx T_e = 100\ \mathrm{eV}$, $\gamma = 5/3$), the sound speed is
\begin{align*}
    c_s = \sqrt{\frac{\gamma k_B T}{m_i}}
        = \sqrt{\frac{5/3 \times 100\ \mathrm{eV} \times 1.6\times10^{-19}\ \mathrm{J\,eV^{-1}}}
        {1.67\times10^{-27}\ \mathrm{kg}}}
        \approx 9.2\times10^4\ \mathrm{m\,s^{-1}},
\end{align*}
giving $v_f \approx 5.6\times10^5\ \mathrm{m\,s^{-1}}$.
For a grid with $\Delta r = \Delta z = \Delta x$, \eqref{eq: CFL} simplifies to
$\Delta t \leq \Delta x / (\sqrt{2}\, v_f)$.
Taking the smallest Polomac length scale $L = 0.04\ \mathrm{m}$ and a modest
$N = 100$ grid points per direction ($\Delta x = 4\times10^{-4}\ \mathrm{m}$),
this gives
\begin{align*}
    \Delta t lesssim \frac{4\times10^{-4}}{\sqrt{2}\times 5.6\times10^5}
        \approx 5\times10^{-10}\ \mathrm{s}.
\end{align*}
This is comfortably smaller than the Alfv\'en timescale
$\tau_A \approx 5.5\times10^{-7}\ \mathrm{s}$, so a simulation spanning
several Alfv\'en times requires of order $10^3$ timesteps.

\subsection{Boundary conditions}
\label{subsec: BCs}

We model the plasma as bounded by perfectly conducting walls, consistent with
the assumption used for the static equilibrium in \autoref{subsec: boundary conditions}.
At a perfectly conducting wall with outward normal $\hat{n}$, the standard
electromagnetic boundary conditions require that no magnetic flux passes through
the wall and no plasma flows through it:
\begin{align}
    \vec{B}_1 \cdot \hat{n} &= 0, \label{eq: BC B} \\
    \vec{V}_1 \cdot \hat{n} &= 0. \label{eq: BC V}
\end{align}
In the $(r,z)$ domain these translate to: at the outer wall ($r = r_\text{wall}$),
$B_r^{(1)} = 0$ and $V_r = 0$; at the top and bottom caps ($z = z_\text{top/bot}$),
$B_z^{(1)} = 0$ and $V_z = 0$. At the inner boundary ($r = 0$ or at the surface of
the central coil), the regularity condition $\partial_r(\cdot)|_{r=0} = 0$ is applied
for scalar quantities, consistent with \autoref{subsec: boundary conditions}.
The tangential components ($V_\varphi$, $B_\varphi^{(1)}$, $V_z$ at radial walls,
$V_r$ at axial walls) are left free, consistent with ideal MHD — a perfectly
conducting wall exerts no tangential force on the plasma.
